\chapter{Repository Map}
\label{app:repo-map}

This appendix summarizes the key repository directories used to implement the CARLA\texorpdfstring{$\rightarrow$}{->}QLabs adaptation and highlights the primary entry points referenced throughout the thesis.

\section{High-level structure}
\begingroup
\singlespacing\small
\begin{verbatim}
Qcar2SimDev/
  core/                  QLabs + QCar2 interface, camera processing, scene spawning
  inference/              Runtime controller, model wrapper, state estimation, control conversion
  data_collection/        Teleoperation + expert data recording scripts
  simlingo/simlingo_training/  Training/eval code + Hydra configuration tree
  config/                 QLabs assets (actors), routes, and scene definitions (JSON)
  results/                Evaluation scripts, baselines, and run logs
  models/simlingo/         Checkpoints + Hydra snapshot used by inference
  report/                 This thesis (LaTeX)
\end{verbatim}
\endgroup

\begin{itemize}[leftmargin=*]
  \item \path{core/}: QLabs/QCar~2 interface, scene loading/spawning, coordinate utilities, camera preprocessing.
  \item \path{inference/}: Runtime controller, model wrapper, state estimator, and control conversion.
  \item \path{data_collection/}: Teleoperation and expert dataset recording tools.
  \item \path{simlingo/simlingo_training/}: Training code (datasets, models, Hydra configs).
  \item \path{config/}: Scene and route JSONs used for both training and evaluation.
  \item \path{results/}: Automated tests, baselines, and evaluation logs.
\end{itemize}

\section{Key entry points}
\begin{itemize}[leftmargin=*]
  \item Inference (runtime control loop): \path{inference/main.py}
  \item Training (PyTorch Lightning + Hydra): \path{simlingo/simlingo_training/train.py}
  \item Data collection (teleop + dataset logging): \path{data_collection/collect_data.py}
\end{itemize}

\section{Module index}
\begin{table}[h]
\centering
\caption{Key implementation modules (non-exhaustive) used in the CARLA\texorpdfstring{$\rightarrow$}{->}QLabs adaptation.}
\label{tab:repo-module-index}
\begin{tabular}{p{0.36\linewidth}p{0.56\linewidth}}
\toprule
Module & Purpose\\
\midrule
\path{core/qcar2_interface.py} & QLabs API integration, QCar~2 spawning/control, camera capture and RGB conversion.\\
\path{core/camera_processor.py} & Image preprocessing to match training-time inputs (e.g., patching/resizing).\\
\path{core/scene_loader.py} & Loads route/scene definitions from \path{config/routes/} and \path{config/scenes/}.\\
\path{core/scene_spawner.py} & Spawns configured actors (vehicles/pedestrians/static props) from \path{config/actors/}.\\
\path{inference/main.py} & Main controller that runs the closed-loop inference+control cycle at runtime.\\
\path{inference/simlingo_model.py} & Loads the fine-tuned checkpoint and formats multimodal prompts/outputs.\\
\path{inference/state_estimator.py} & Estimates ego state (pose/velocity) used by the controller.\\
\path{inference/control_converter.py} & Converts model outputs into low-level QCar~2 steering/throttle commands.\\
\path{data_collection/data_recorder.py} & Persists expert demonstrations (images + measurements) into SimLingo dataset format.\\
\path{results/test_simlingo_roundabout.py} & Scenario-based evaluation harness and safety/route metrics.\\
\path{results/test_acc_baseline.py} & Baseline controller (ACC-style) used for comparison.\\
\bottomrule
\end{tabular}
\end{table}

\section{Configuration and assets}
\begin{itemize}[leftmargin=*]
  \item Scenes (which route + actors are active): \path{config/scenes/training/} and \path{config/scenes/testing/}.
  \item Routes (waypoints/spawn definitions): \path{config/routes/*.json}.
  \item Actor templates (vehicles/pedestrians/static props): \path{config/actors/}.
  \item Training configs (Hydra): \path{simlingo/simlingo_training/config/}.
  \item Run-time model config snapshot used by inference: \path{models/simlingo/.hydra/config.yaml}.
\end{itemize}
